\hwpchapter{4}{현장 활용방안}{10}

\hwpsection{4.1 예측 결과의 현장 활용}

C-BAT는 매일 00:00에 당일의 시간별 생산계획과 전날까지 관측된 전력으로 하루 96슬롯의 예측 경로(사후 표본)를 만들고, 이 경로에서 점검에 필요한 값을 뽑아 일일 보고로 낸다(표~\ref{tab:ch4-report}). 보고는 설비 담당자가 어느 시간대와 어느 날을 먼저 점검할지 정하는 데 쓰며, 경보만으로 설비를 자동 제어하거나 정지시키지 않는다.

\begin{table}[htbp]
\hwptablefont
\setlength{\tabcolsep}{4pt}
\caption{00:00 일일 보고의 항목과 읽는 방법.}\label{tab:ch4-report}
\begin{tabular}{L{2.5cm}L{4.4cm}L{3.8cm}L{4.4cm}}
\toprule
항목 & 내용 & 현장에서 쓰는 방법 & 읽을 때의 한계 \\
\midrule
분위수 곡선 & 슬롯별 예측 분위수 19개(5\%p 간격, 5--95\%) & 부하가 높아질 시간대와 불확실한 시간대를 보고 그 시간의 생산량과 설비 동시 가동을 점검 & 슬롯 90\% 구간 적중률이 개발 구간에서 0.990으로 명목값보다 높아 구간이 넓다 \\
일 피크 분포 & 경로별 하루 최댓값의 분포(중앙값, 평균) & 그날 일 피크의 수준과 불확실성 확인 & 일 피크 90\% 구간 적중률 0.948, 평균 폭 49.18 kW(개발 구간) \\
피크 시각 & 경로별 피크 시각의 최빈값 & 점검할 시간대의 후보 & 관측 피크 시각의 $\pm$30분 안에 든 비율이 0.203(개발 구간)이라 시각을 단정하지 않는다 \\
피크 기준 초과 확률 & 일 피크가 C50·C75·C90을 넘는 경로의 비율(보정하지 않음) & 점검할 날의 우선순위(3장의 피크 기준 초과 경보) & 7월 종일 가동일의 미탐지와 평범한 가동일의 오경보가 남는다(3장) \\
기준선 $B_d$ & 예측일 00:00에 알려진 요금 적용 최대(4.2절) & 월 최대 갱신 여부를 가르는 기준 & 이전 자료만 쓰며, 갱신이 일어나면 다음 날부터 올라간다 \\
$P_d$, $E_d$ & 경로 최댓값이 $B_d$를 넘을 확률 $P_d$와 넘은 만큼의 기본요금 기댓값 $E_d$(4.2절) & 월 최대 갱신 경보 & 갱신일이 1일뿐이어서 사례로만 읽는다 \\
\bottomrule
\end{tabular}
\end{table}

이 값들은 시간별 생산계획을 미리 안다는 가정(1.4절)에서 나왔다. 계획 입력에 오차를 넣은 실험에서 C-BAT의 슬롯 MAE는 생산량 오차(로그 표준편차 0.2)에서 0.04\% 늘었고 계획 전체를 1시간 이르게 넣으면 3.3\% 늘었다. 현장에서는 생산량보다 가동 시작·종료 시각의 계획 정확도가 더 중요하다. 계획 취소나 설비 고장처럼 가동 여부 자체가 바뀌는 경우는 시험하지 않았다.

\hwpsection{4.2 월 최대 갱신 경보 규칙과 백테스트}

기본요금은 매달 요금 시간대의 최대수요와 래칫 바닥 가운데 큰 값(청구 수요)에 1 kW당 8,320원(2021년 산업용(을) 요금표)을 곱해 정해진다. 요금 시간대는 일요일과 공휴일을 뺀 09--23시이고, 래칫 바닥은 지난 1년 동안 래칫 적용월(12·1·2·7·8·9월)의 최대수요에서 온다. 따라서 기본요금을 움직이는 날은 그 달의 최대를 새로 쓰는 날뿐이다. 이 장에서는 이 날을 월 최대 갱신일, 그 위험을 당일 00:00에 알리는 경보를 월 최대 갱신 경보라 한다. 이 경보는 일 피크가 C50·C75·C90을 넘을 확률로 내는 3장의 피크 기준 초과 경보와 기준이 다르다.

경보 규칙은 백테스트 결과를 보기 전에 문서로 정해 두었고, 예측일 $d$의 00:00에 다음 순서로 계산한다.
\begin{hwplist}
\item 기준선: 이번 달 1일부터 전날까지의 요금 시간대 관측 최대와, 이번 달 이전 1년 안 래칫 적용월의 요금 시간대 관측 최대 가운데 큰 값을 기준선 $B_d$로 둔다. $d$일 이전 관측만 쓴다.
\item 경로 최댓값: C-BAT의 사후 경로 $s=1,\dots,S$마다 요금 시간대 슬롯의 최댓값 $M_d^{(s)}$를 구한다.
\item 위험: 갱신 확률 $P_d$와 기대 초과 기본요금 $E_d$를 식~\eqref{eq:ch4-pe}로 구한다. 시드 3개(0·1·2)의 값을 평균한다.
\item 경보: $E_d$가 점검 비용 3만 원(담당자 1--2시간 가정) 이상이면 경보한다(주 규칙). 1만 원, 10만 원, $P_d \ge 0.5$는 민감도로 본다.
\end{hwplist}
\begin{equation}
P_d = \frac{1}{S}\sum_{s=1}^{S} \mathbf{1}\bigl[M_d^{(s)} > B_d\bigr], \qquad
E_d = 8{,}320 \times \frac{1}{S}\sum_{s=1}^{S} \max\bigl(0,\ M_d^{(s)} - B_d\bigr).
\label{eq:ch4-pe}
\end{equation}
위험은 Platt 보정을 거치지 않은 경로 비율을 그대로 쓴다. 보정에 쓰는 최근 7일 안에 C90 초과일이 없으면 보정식이 절편만 추정되어 모든 날의 위험이 같은 값이 되기 때문이다. 개발 구간에서도 C90의 AUC는 보정 전 0.926에서 보정 뒤 0.788로 낮아졌다.

백테스트는 시뮬레이션이 아니라 과거 각 날에 규칙을 적용하고 실제 관측으로 채점한 것이다. 대상은 개발 구간(2021년 7월 7일--8월 31일) 가운데 요금 시간대 관측이 있는 44일이며, 이 중 월 최대 갱신일은 7월 19일 하루였다. 그날 요금 시간대 최대는 222 kW로 전날까지의 기준선 206 kW를 16 kW 넘었고, 기본요금 8,320원/kW를 적용하면 그 달 기준 133,120원이다. 주 규칙은 이날을 잡았고 44일 동안 10번 경보하였다(표~\ref{tab:ch4-alarm}). 기준일 곡선을 그대로 쓰는 B0\_kind는 이날을 놓쳤다.

\begin{table}[htbp]
\centering
\hwptablefont
\caption{월 최대 갱신 경보 백테스트(개발 구간 중 요금 시간대 관측이 있는 44일, 갱신일 1일).}\label{tab:ch4-alarm}
\begin{tabular}{lcrrr}
\toprule
규칙 & 갱신일 포착 & 경보 수 & 오경보 & 금액 상한(원) \\
\midrule
$E_d \ge$ 1만 원 & 1/1 & 23 & 22 & 133,120 \\
$E_d \ge$ 3만 원(주 규칙) & 1/1 & 10 & 9 & 133,120 \\
$E_d \ge$ 10만 원 & 1/1 & 3 & 2 & 133,120 \\
$P_d \ge 0.5$ & 1/1 & 3 & 2 & 133,120 \\
C-BAT 점예측 $> B_d$ & 1/1 & 3 & 2 & 133,120 \\
B0\_kind 기준일 $> B_d$ & 0/1 & 0 & 0 & 0 \\
매일 경보 & 1/1 & 44 & 43 & 133,120 \\
\bottomrule
\end{tabular}
\end{table}
\hwpnote{금액 상한은 갱신일의 초과분을 완벽한 조치로 모두 막았다고 가정한 그 달 기준 값이며 실제 절감액이 아니다. C-BAT 점예측은 경로 최댓값의 중앙값이다. 2020년 12월 자료가 없어 기준선에는 그 달의 최대가 들어 있지 않다.}

7월 19일에 C-BAT는 갱신 확률 $P_d=0.58$, 기대 초과 기본요금 $E_d=122{,}919$원, 경로 최댓값의 중앙값 211.19 kW를 냈다. B0\_kind의 기준일 곡선은 이날 요금 시간대 최대를 195 kW로 두어 기준선에 못 미쳤다. 점검 비용 기준을 10만 원으로 올리면 경보는 3번으로 줄지만 갱신일은 여전히 잡힌다. 오경보 9건 가운데 6건은 7월 7--16일에 몰려 있다. 이 기간 관측 피크는 187--201 kW였으나 C-BAT 경로 최댓값의 중앙값은 201.56--214.65 kW로 기준선 206 kW 부근이었다. 기준선이 평소 가동일 피크와 가까웠고, 평범한 가동일의 피크를 높게 예측하는 C-BAT의 경향(3장)이 겹친 결과이다. 갱신일이 1일뿐이므로 이 결과는 성능 추정이 아니라 사례이다.

\hwpsection{4.3 경보일 점검 항목과 금액 해석}

경보일에는 다음 항목을 점검한다. 항목은 관측 조건에서 정한 것이며 저감 효과는 검증하지 않았다. 최종 적용은 설비·작업 순서·납기를 아는 담당자가 정한다.
\begin{hwplist}
\item 초반 부하: 유사일 비교에서 초반 두 시간의 생산 비중이 10\%p 높은 날은 요금 시간대 피크가 9.3 kW 높았다. 실제 설비 가동 기록으로 동시 기동 여부를 확인하고, 작업 제약이 허용하면 기동 순서를 분산하는 조치를 시험한다.
\item 08시·11시 부하 집중: 개발 구간 고피크일 13일 가운데 10일(08시 6일, 11시 4일)이 두 시각 중 하나에 피크를 기록하였다(1.2절). 이 집계만으로 동시 기동이 원인이라고 단정하지 않는다. 08시는 요금 시간대 밖이라 기본요금 산정에 들어가지 않고, 11시는 여름철 최대부하 시간대로 산정에 들어간다.
\item 기준선 근접 시 비필수 부하 지연: 당일 요금 시간대 부하가 $B_d$에 가까워지면 비필수 부하를 뒤로 미룬다.
\end{hwplist}

초반 부하 항목의 방향은 모형과 무관한 유사일 비교에서 얻었다. 2021년 9월 이전 가동일 89일에서 가동유형과 날짜유형이 같고 일 생산량 차이가 10\% 이내인 날로 쌍을 만들고, 시작 시각이나 초반 생산 비중만 다른 쌍의 실측 요금 시간대 피크를 비교하였다. 시작 시각이 1시간 이상 다른 206개 쌍에서는 늦게 시작한 날의 피크가 평균 11.1 kW 높았다(95\% 구간 $[5.2,\ 17.4]$). 피크 차이의 합을 시작 지연 시간의 합으로 나눈 값은 1시간당 7.4 kW였다(95\% 구간 $[3.4,\ 12.0]$). 그러나 쌍별 피크 차이의 중앙값은 2.0 kW로, 큰 차이는 일부 쌍에 몰려 있다. 시작 시각이 같고 초반 두 시간의 생산 비중이 5\%p 이상 다른 쌍은 25개뿐이었고 10\%p당 9.3 kW였다(95\% 구간 $[0.6,\ 20.2]$). 한 날이 여러 쌍에 들어가 구간은 실제보다 좁을 수 있고, 시작 시각은 수주량이나 설비 상태와 함께 움직일 수 있다. 따라서 이 값은 관측된 연관성이며 조치의 효과 크기로 쓰지 않는다.

금액은 갱신일을 완벽하게 막았다고 가정한 상한과 점검 비용을 나란히 놓고 읽는다. 주 규칙의 경보일 집합을 $\mathcal A$, 월 최대 갱신일 집합을 $\mathcal E$, 갱신일 $d$의 요금 시간대 관측 최대를 $M_d^{\mathrm{obs}}$라 하면 그 달 기준 계산은 다음과 같다.
\begin{align}
U_{\mathrm{month}} &= 8{,}320\sum_{d\in\mathcal A\cap\mathcal E}
\bigl(M_d^{\mathrm{obs}}-B_d\bigr)^+ = 133{,}120\ \text{원},\notag\\
K_{\mathrm{inspect}} &= 30{,}000\times|\mathcal A| = 30{,}000\times 10 = 300{,}000\ \text{원},\label{eq:ch4-cost}\\
U_{\mathrm{month}} - K_{\mathrm{inspect}} &= -166{,}880\ \text{원}.\notag
\end{align}
이는 점검 비용을 10번의 경보일마다 실제로 지출하고, 포착한 갱신일의 초과분을 전부 막는 경우의 계산이다. 당월 계산만으로는 순편익이 양수라는 근거가 없다. 7월은 래칫 적용월이어서 7월 19일의 초과분은 이후 달의 청구 수요에도 영향을 줄 수 있다. 실제로 8월 요금 시간대 최대는 207 kW였지만 8월 청구 수요는 7월이 만든 222 kW였다. 이 이후 달 영향은 이후 달의 실제 피크에 따라 달라지므로 금액에 넣지 않았다. 점검 시간, 조치 비용, 후속 월 청구 수요는 현장에서 기록해 순편익을 평가해야 한다.
\hwpnote{133,120원은 초과분 16 kW 전체를 줄이는 느슨한 상한이다. 7월의 다음 최대는 7월 21일의 214 kW여서 7월 19일만 낮추면 7월 청구 수요는 214 kW까지 내려가며, 이 경우 감소분은 8 kW $\times$ 8,320원 = 66,560원이다.}

\hwpsection{4.4 적용 절차와 효과 확인}

모형이 계산하는 시간대 이동의 반사실 절감액은 현장 금액으로 쓰지 않았다. 시간대 이동은 가동 시작 시각을 앞이나 뒤로 옮기는 등 하루 총생산량은 두고 생산계획만 바꾸는 변경이다. C-BAT에서 생산계획만 바꾸어 곡선을 다시 계산하면 절감액을 낼 수 있으나, 그 금액이 관측과 맞는지 먼저 확인하였다. 4.3절의 유사일 쌍 231개(206개 + 25개)에서 두 날의 실측 차이와, 두 날의 계획을 서로 바꿔 넣은 모형 차이를 같은 요금 달력으로 비교하였다. 전력량(kWh)은 슬롯 전력(kW)에 0.25시간을 곱해 합한 값이고, 전력량요금은 여기에 시간대별 단가를 곱한 요금이다. 결과는 표~\ref{tab:ch4-counter}와 같다.

\begin{table}[htbp]
\centering
\hwptablefont
\setlength{\tabcolsep}{4pt}
\caption{유사일 231쌍에서 모형 차이와 실측 차이의 비교(늦게 시작하거나 초반 생산 비중이 높은 날 $-$ 비교 날).}\label{tab:ch4-counter}
\begin{tabular}{lrrrl}
\toprule
대상 & 실측 평균 & 모형 평균 & 상관 & 기울기 [날짜 단위 95\% 구간] \\
\midrule
전력량요금 (원/일) & $+7{,}359$ & $-464$ & $-0.01$ & $-1.39$ [$-15.35$, 8.25] \\
요금 시간대 피크 (kW) & $+11.0$ & $+0.2$ & 0.18 & 9.51 [$-10.39$, 19.09] \\
\bottomrule
\end{tabular}
\end{table}
\hwpnote{기울기는 실측 차이를 모형 차이에 원점 회귀한 값으로, 모형 값에 몇 배를 곱해야 실측과 맞는지를 나타낸다. 구간은 날짜 단위 cluster 부트스트랩이며 쌍 단위 구간은 각각 $[-5.06,\ 2.13]$, $[0.07,\ 14.40]$이다. C-BAT는 9월 이전 전체 자료로 한 번 적합하였으므로 쌍의 날짜가 학습에 들어 있는 표본 내 검증이다.}

전력량요금은 상관이 0에 가깝고 평균 방향도 반대여서, 실측에서는 늦게 시작하거나 초반 생산 비중이 높은 날이 하루 7,359원 비쌌지만 모형은 464원 싸다고 계산하였다. 요금 시간대 피크는 모형의 평균 반응이 0.2 kW로 실측 차이 11.0 kW에 크게 못 미쳤고, 날짜 단위 구간이 0을 포함하였다. 실측 쌍은 계획 외의 조건도 다를 수 있어 인과 기준값은 아니다. 그래서 모형이 틀렸다고 결론짓지 않고, 모형 절감액을 관측으로 뒷받침할 수 없다는 데서 멈춘다.

실제 기본요금 규칙을 적용해도 같은 결론이다. 모형이 요금 시간대 피크 감소를 이유로 시간대 이동을 제안한 날은 C-BAT에서 7월 0일, 8월 1일이었다. 7월 청구 수요 222 kW를 정한 7월 19일은 이 날들에 들지 않아 7월 기본요금은 변하지 않는다. 8월은 7월이 만든 래칫 222 kW가 청구 수요를 정해 8월 피크를 얼마나 낮추든 기본요금이 같다. 그래서 두 달의 기본요금 변화는 모두 0원이었고, 모형의 피크 감소를 10배로 키워도 0원이었다. 이에 따라 이 보고서의 금액 주장은 관측된 요금 구조, 즉 갱신일 하루가 만든 금액 상한까지로 제한한다.

조치의 효과는 현장 운영으로 검증한다. 작업 제약이 허용하는 경보일을 사전에 정한 무작위 배정으로 조치일과 비조치일에 나누고, 가동유형과 계획 생산량을 함께 기록하여 두 집단의 요금 시간대 피크를 비교한다. 갱신일이 드물어 비교 지표는 갱신 여부보다 경보일의 요금 시간대 피크를 우선한다. 조치하기 쉬운 날만 고르지 않도록 배정 규칙과 미이행 사유를 남기고, 실제 점검 비용과 조치 비용을 포함한 순편익도 기록한다. 도입은 다음 세 단계로 진행한다.
\begin{hwplist}
\item 1단계(매일 00:00): 당일 생산계획으로 C-BAT를 실행해 분위수 곡선, 일 피크 분포, 피크 시각, $B_d$, $P_d$, $E_d$를 담은 일일 보고를 만든다.
\item 2단계(경보일): 담당자가 보고를 확인하고 설비·작업 순서·납기를 고려하여 4.3절의 점검 항목 가운데 적용할 것을 정한다. 무작위 배정 결과와 적용 여부, 내용을 기록한다.
\item 3단계(매월 말): 경보 수, 오경보, 갱신일 포착 여부, 조치일과 비조치일의 요금 시간대 피크를 검토하고 점검 비용 기준을 조정한다.
\end{hwplist}
\hwpnote{이 장의 경보 백테스트는 개발 구간 중 요금 시간대 관측이 있는 44일에 한정하며, 합산 평가(67일)의 피크 예측 성능과 구별한다. 월 최대 갱신 사례는 1건이어서 경보 규칙의 성능을 추정하지 않았다.}
